%
% crosshost_hubid_preemption.tex
%
% W12 (lux-833p), Invariant 2 of the cross-host transport design (DES-090):
% "at most one live connection per HubId."  Models the single-owner
% preemption discipline of DES-068 as re-keyed by W11 (lux-u9gb) onto the
% connecting Hub's HubId, generalized from one same-host Hub to N
% concurrently-reconnecting Hubs.
%
% Models the code in:
%   src/punt_lux/display/client_registry.py   -- ClientRegistry.hub_fd_for
%   src/punt_lux/display/hub_reconciliation.py -- HubReconciliation._preempt_stale_hub
% and the invariant stated in
%   docs/architecture/system.tex "Cross-Host Transport and Trust",
%   subsection "Invariants", item 2.
%
% This document supersedes the name-keyed preemption clauses of
% hub_display_reconciliation.tex (its own SUPERSEDED banner points here):
% HI4 there ("preemption scopes to the shared name") described the pre-W11
% code and is corrected below.
%
% Type-check:  fuzz docs/crosshost_hubid_preemption.tex
% Model-check (goal predicates in the Verification section):
%   probcli docs/crosshost_hubid_preemption.tex -model_check -p DEFAULT_SETSIZE 3
%
% Formatting follows the fuzz house rules: \quad~ for continuation indent,
% schema lines under 80 columns, predicates broken at \land/\lor/\implies.
%

\documentclass[a4paper,10pt,fleqn]{article}
\usepackage[margin=1in]{geometry}
\usepackage{fuzz}
\usepackage[colorlinks=true,linkcolor=blue,citecolor=blue,urlcolor=blue]{hyperref}

\begin{document}

\title{Cross-Host Single-Owner Preemption: a Z Specification}
\author{W12 formal model of DES-090 Invariant 2 (at most one live
connection per \texttt{HubId})}
\date{September 2026}
\maketitle

% ============================================================
\section{Introduction}
% ============================================================

The Display may aggregate Hub connections from several machines at once.
Each such connection carries, in its \texttt{ConnectMessage}, two
independent identifiers: a \texttt{HubId} --- a per-process identity, the
peer's hostname (cross-host, certificate-verified) paired with its
process id --- and a \texttt{name}, a human label that many Hub processes
may legitimately share (every production Hub today declares the identical
hardcoded name).  The two are not the same kind of thing, and the whole
point of this model is that single-owner preemption must key on the
former, never the latter.

DES-068 established, for the then-single same-host Hub, that a fresh
\texttt{kind="hub"} identify forcibly disconnects any predecessor
declaring the \emph{same identity} before recording itself --- so a dead
process's stale connection cannot linger and swallow clicks.  With one
Hub, ``same identity'' was unambiguous.  Cross-host makes two
independently-legitimate, simultaneously-live Hubs the ordinary case, and
the question ``which predecessor does a new connection preempt?'' now has
a right answer and a wrong one:

\begin{description}
  \item[right --- key on \texttt{HubId}.] A reconnect under the same
    \texttt{HubId} preempts its own stale predecessor; two connections
    with \emph{different} \texttt{HubId}s never preempt each other, even
    when they share a \texttt{name}.  This is the W11 (\texttt{lux-u9gb})
    change: \texttt{ClientRegistry.hub\_fd\_for} keys the lookup on
    \texttt{HubId}.
  \item[wrong --- key on \texttt{name}.] The pre-W11 code keyed the
    lookup on \texttt{ConnectMessage.name}.  Two Hubs sharing a name on
    different machines would then evict one another (a live, unrelated
    Hub wrongly disconnected), while two connections sharing a
    \texttt{HubId} but declaring different names would \emph{fail} to
    collapse --- both left live for one identity.
\end{description}

Two properties must hold across every interleaving of $N$ Hubs
connecting, reconnecting, and disconnecting:

\begin{description}
  \item[I1 --- at most one live owner per \texttt{HubId}.] No reachable
    state holds two distinct live connections declaring the same
    \texttt{HubId}.  This is Invariant 2 of the design, stated exactly.
  \item[I2 --- deadlock freedom.] No interleaving traps the model; some
    operation is always enabled.
\end{description}

Alongside these, two \emph{witnesses} must be reachable, or the
invariants would hold only because the mechanism they guard is dead
(Section~\ref{sec:verify}): different-\texttt{HubId} connections
\emph{coexisting} (preemption is not over-eager), and a same-name pair of
them coexisting in particular (the coexistence the wrong, name-keyed
design destroys).

% ============================================================
\section{Basic Types}
% ============================================================

\begin{zed}
  [CONN, HID, NAME]
\end{zed}

\texttt{CONN} is the set of Hub-Display connections (file descriptors)
the Display can hold open; three suffice to model two connections
contending over one identity plus a third unrelated one.  \texttt{HID} is
the set of \texttt{HubId} values a connection may declare --- two suffice
to exhibit both a same-identity reconnect and two distinct coexisting
identities.  \texttt{NAME} is the set of declared human labels; two
suffice to exhibit the (same \texttt{HubId}, different \texttt{name}) pair
the name-keyed variant mis-handles and the (different \texttt{HubId},
same \texttt{name}) pair it wrongly evicts.  All three carriers are
bounded by ProB's \texttt{DEFAULT\_SETSIZE}.

% ============================================================
\section{State}
% ============================================================

The state is flat: one schema.  A connection is ``live'' once it has
identified as \texttt{kind="hub"} and its socket is still open --- the
state in which it owns content and services clicks.  \texttt{hubOf} and
\texttt{nameOf} record what each live connection declared.

\begin{schema}{Owners}
  live : \power CONN \\
  hubOf : CONN \pfun HID \\
  nameOf : CONN \pfun NAME
\where
  \dom hubOf = live \\
  \dom nameOf = live
\end{schema}

The two domain equalities are structural, not the safety property: a live
connection has, by construction, declared exactly one \texttt{HubId} and
one \texttt{name} (both arrive in its single \texttt{ConnectMessage}), and
a departed one has neither.  ``At most one live per \texttt{HubId}'' ---
the injectivity of \texttt{hubOf} on \texttt{live} --- is I1, deliberately
\emph{not} written here: placing it in the state schema would enforce it
as a guard on every successor and silently disable any operation that
could break it, masking the very defect the name-keyed variant
(Section~\ref{sec:fidelity}) must be able to reach.

% ============================================================
\section{Initialization}
% ============================================================

\begin{schema}{Init}
  Owners'
\where
  live' = \emptyset \\
  hubOf' = \emptyset \\
  nameOf' = \emptyset
\end{schema}

A single initial state: no connection is live, and nothing has been
declared --- the cold start before any Hub has connected.

% ============================================================
\section{Operations}
% ============================================================

\subsection{Identify --- single-owner preemption, keyed on \texttt{HubId}}

A fresh connection \texttt{c?} identifies as \texttt{kind="hub"} declaring
\texttt{HubId} \texttt{h?} and name \texttt{n?}.  \texttt{c?} must not
already be live: a reconnect always arrives on a new file descriptor, so
the predecessor is a \emph{different} \texttt{CONN} value, and idempotent
re-identification of an already-live fd is not a distinct transition here.

The preempted set is exactly the live connections declaring the same
\texttt{HubId} \texttt{h?} --- \texttt{hub\_fd\_for(h?)} in the code,
lifted from ``the one such fd'' to ``every such fd'' so the model can
never assume uniqueness it is trying to prove.  Each is dropped
(\texttt{remove\_client}); \texttt{c?} is then recorded as the sole live
owner of \texttt{h?}.

\begin{schema}{Identify}
  \Delta Owners \\
  c? : CONN \\
  h? : HID \\
  n? : NAME
\where
  c? \notin live \\
  live' = (live \setminus \{ c : live | hubOf~c = h? \}) \cup \{ c? \} \\
  hubOf' = (\{ c : live | hubOf~c = h? \} \ndres hubOf) \oplus \{ c? \mapsto h? \} \\
  nameOf' = (\{ c : live | hubOf~c = h? \} \ndres nameOf) \oplus \{ c? \mapsto n? \}
\end{schema}

\subsection{Disconnect --- an ordinary drop}

A live connection's socket closes on its own --- a network blip, a
graceful shutdown, or a process exit.  \texttt{c?} leaves \texttt{live}
and its declarations are forgotten.  No new mechanism: a dropped
cross-host connection is, to the socket listener, an ordinary fd closing,
handled by the existing per-fd reaping the same-host case already uses.

\begin{schema}{Disconnect}
  \Delta Owners \\
  c? : CONN
\where
  c? \in live \\
  live' = live \setminus \{ c? \} \\
  hubOf' = \{ c? \} \ndres hubOf \\
  nameOf' = \{ c? \} \ndres nameOf
\end{schema}

% ============================================================
\section{Invariants}
\label{sec:inv}
% ============================================================

\paragraph{Invariant 1 --- at most one live owner per \texttt{HubId}.}
\texttt{hubOf} is injective on the live set:
\[
  \forall c_1, c_2 : live @ \\
  \quad~ hubOf~c_1 = hubOf~c_2 \implies c_1 = c_2.
\]
\texttt{Identify} removes \emph{every} live connection declaring
\texttt{h?} before adding \texttt{c? \mapsto h?}, so no two live
connections can share a \texttt{HubId}; \texttt{Disconnect} only shrinks
\texttt{live}, which cannot create a collision.  The check
(Section~\ref{sec:verify}) asks whether the negation --- two distinct live
connections with equal \texttt{hubOf} --- is reachable.

\paragraph{Invariant 2 --- deadlock freedom.}
\texttt{Disconnect} is enabled whenever \texttt{live} is non-empty, and
\texttt{Identify} is enabled for any \texttt{c? \notin live} whenever one
exists.  In the sole state where \texttt{Identify} has no fresh argument
--- \texttt{live = CONN} --- \texttt{Disconnect} is enabled.  So some
operation is always enabled.

% ============================================================
\section{Verification}
\label{sec:verify}
% ============================================================

Type-checking is a \texttt{fuzz} pass:
\begin{verbatim}
  fuzz docs/crosshost_hubid_preemption.tex
\end{verbatim}

I1 is a state property, checked by asking ProB whether its negation is
reachable.  A goal reported \emph{not found} means the invariant holds on
every reachable state.  It returns \emph{goal not found} against
\texttt{DEFAULT\_SETSIZE 3}:
\begin{verbatim}
  # Invariant 1 (goal NOT found for the correct spec)
  probcli docs/crosshost_hubid_preemption.tex -model_check -nodead \
    -goal "#(c1,c2).(c1 : live & c2 : live & not(c1 = c2) &
                     hubOf(c1) = hubOf(c2))" \
    -p DEFAULT_SETSIZE 3
\end{verbatim}

I2 is the full-state-space deadlock check.  Every state has an enabled
operation (Section~\ref{sec:inv}), so any deadlock reported would be a
model defect:
\begin{verbatim}
  # Invariant 2 / deadlock freedom (no deadlock, all states visited)
  probcli docs/crosshost_hubid_preemption.tex -model_check \
    -p DEFAULT_SETSIZE 3
\end{verbatim}

The witnesses must each be \emph{found}, or the invariant above would
hold only because the mechanism it guards is unreachable.  Both return
\emph{goal found} for the correct spec:
\begin{verbatim}
  # W1: two different-HubId connections coexist -- preemption is not
  #     over-eager (goal FOUND)
  probcli docs/crosshost_hubid_preemption.tex -model_check -nodead \
    -goal "#(c1,c2).(c1 : live & c2 : live & not(c1 = c2) &
                     not(hubOf(c1) = hubOf(c2)))" \
    -p DEFAULT_SETSIZE 3

  # W2: two different-HubId connections that share a NAME coexist --
  #     the coexistence the name-keyed variant destroys (goal FOUND)
  probcli docs/crosshost_hubid_preemption.tex -model_check -nodead \
    -goal "#(c1,c2).(c1 : live & c2 : live & not(c1 = c2) &
                     not(hubOf(c1) = hubOf(c2)) &
                     nameOf(c1) = nameOf(c2))" \
    -p DEFAULT_SETSIZE 3
\end{verbatim}

% ============================================================
\section{Fidelity: reproducing the shipped defect}
\label{sec:fidelity}
% ============================================================

A model that cannot reproduce the known defect proves nothing.  The buggy
variant, kept as
\texttt{docs/crosshost\_hubid\_preemption\_name\_keyed\_buggy.tex}, is the
pre-W11 design: preemption keys on \texttt{name}, not \texttt{HubId}.
\texttt{Identify}'s preempted set becomes
\begin{verbatim}
  % correct (this spec):
  preempted == { c : live | hubOf(c) = h? }

  % buggy (name-keyed variant):
  preempted == { c : live | nameOf(c) = n? }
\end{verbatim}
and nothing else changes.  That single substitution reproduces both
halves of the defect W11 closes.

\paragraph{The double-owner trace (I1 violated).}
\begin{enumerate}
  \item \texttt{Identify}(c1, h1, n1): \texttt{live = \{c1\}},
    \texttt{hubOf = \{c1 \mapsto h1\}}, \texttt{nameOf = \{c1 \mapsto n1\}}.
  \item \texttt{Identify}(c2, h1, n2): the same \texttt{HubId} \texttt{h1},
    but a \emph{different} name \texttt{n2}.  Name-keyed preemption looks
    for a live connection named \texttt{n2}; \texttt{c1} is named
    \texttt{n1}, so nothing is preempted.  \texttt{c2} joins:
    \texttt{live = \{c1, c2\}} with \texttt{hubOf(c1) = hubOf(c2) = h1}.
\end{enumerate}
Two distinct live connections now own the identical \texttt{HubId} ---
I1's negation, reachable.  I1's goal returns \emph{found} against this
variant and \emph{not found} against the correct spec.  This is the
``double-owner state'' the design's fidelity requirement names: the same
\texttt{HubId} claimed by two live connections at once, which
\texttt{HubId}-keyed preemption excludes because step 2 would then find
\texttt{c1} (same \texttt{HubId}) and evict it.

\paragraph{The wrong-eviction trace (W2 destroyed).}
The same substitution also breaks the coexistence witness W2.  Under the
name-keyed variant, any second \texttt{Identify} sharing a \texttt{name}
with a live connection evicts it, so two live connections can never share
a name --- even when their \texttt{HubId}s differ and both are legitimate
Hubs on different machines.  W2's goal, \emph{found} for the correct spec,
returns \emph{not found} against this variant.  That is the live,
unrelated Hub wrongly disconnected because it happened to share the
hardcoded name --- the second half of what W11 fixes.

Both discriminators come from the one key choice, which is why W11 is a
single change, not two, and why this one buggy variant is sufficient
fidelity for Invariant 2.

\paragraph{Errata for \texttt{hub\_display\_reconciliation.tex}.}
That document's HI4 (\texttt{hub\_display\_reconciliation\_coverage.md}:
``preemption scopes to the shared name'') and its \texttt{hubOwner}
identity, modeled as a bare \texttt{CONNECTION\_ID} with no
name/\texttt{HubId} distinction, describe the pre-W11 code.  The corrected
statement is I1 above: preemption scopes to the shared \texttt{HubId}; two
connections sharing a name but declaring distinct \texttt{HubId}s coexist,
and only a same-\texttt{HubId} arrival preempts.  This model, not that
one, is authoritative for Invariant 2.

% ============================================================
\section{Test Partitions}
\label{sec:partitions}
% ============================================================

Derived from the operations and invariants above, for audit against the
test suite (see \texttt{docs/crosshost\_transport\_coverage.md}):

\begin{description}
  \item[Q1 --- a same-\texttt{HubId} reconnect preempts its predecessor.]
    A second \texttt{Identify} with the \texttt{HubId} of a live
    connection evicts it; exactly one live owner remains.
  \item[Q2 --- a different-\texttt{HubId} identify preempts nothing.] Two
    connections with distinct \texttt{HubId}s coexist (witness W1).
  \item[Q3 --- a same-name, different-\texttt{HubId} pair coexists.] A
    shared \texttt{name} does not trigger preemption (witness W2); the
    regression against name-keying.
  \item[Q4 --- a same-\texttt{HubId}, different-name reconnect still
    preempts.] The declared name plays no role: a reconnect under the
    same \texttt{HubId} but a changed name still evicts its predecessor.
  \item[Q5 --- an ordinary disconnect releases the identity.] After
    \texttt{Disconnect}, the freed \texttt{HubId} may be claimed again
    with no residual owner.
  \item[Q6 --- N-way interleaving preserves one-per-\texttt{HubId}.] Any
    interleaving of several Hubs connecting, reconnecting, and
    disconnecting leaves \texttt{hubOf} injective on \texttt{live} (I1,
    model-checked exhaustively).
\end{description}

\end{document}
